\section{Related Work}
\label{sec:related_work}

\subsection{Text-to-SQL Methods}
\label{subsec:text2sql}

Text-to-SQL research has evolved through three paradigms. 
Early fine-tuning approaches~\cite{yu2018spider} trained encoder-decoder models, achieving high accuracy on Spider but suffering from poor cross-database generalization. 
The in-context learning (ICL) paradigm eliminated fine-tuning, with methods like \textsc{Act-SQL}~\cite{actsql2023} and \textsc{Dail-SQL}~\cite{gao2024dail} achieving strong cross-database transfer through skeleton retrieval and example selection. 
The agent paradigm embeds SQL generation in multi-step pipelines.
\textsc{Din-SQL}~\cite{pourreza2023din} introduced task decomposition and self-correction. \textsc{Mac-SQL}~\cite{wang2025macsql} proposed a multi-agent collaborative framework with a planner, selector, and refiner. 
\textsc{Chase-SQL}~\cite{chasesql2024} further extended this with multi-path reasoning and preference optimization. \textsc{Chess}~\cite{talaei2024chess} focused on contextual harness for efficient SQL synthesis.
\textsc{ReFoRCE}~\cite{reforce2025} tops the Spider~2.0 leaderboard by combining self-refinement, consensus enforcement, and column exploration. Despite these advances, all existing agent-based methods treat each query independently---they cannot retain knowledge from past errors.
The enterprise benchmark \textsc{Beaver}~\cite{chen2024beaver} confirms that SQL failures are concentrated in a few recurring categories (schema linking, analytical errors), highlighting the need for cross-query memory. 
Our work addresses this gap by introducing a persistent skill memory that accumulates and generalizes from execution experience.

\subsection{Case-Based Reasoning for Text-to-SQL}
\label{subsec:cbr}

CBR's four-stage cycle (Retrieve-Reuse-Revise-Retain) maps naturally to LLM-based SQL generation. 
Static Retrieve-Reuse methods use fixed case bases---e.g., skeleton matching in \textsc{Act-SQL}~\cite{actsql2023} or de-semanticized retrieval in~\cite{guo2023cbr}---but lack a Retain stage, so the case base cannot grow. Dynamic Retain methods have emerged recently. 
\textsc{ReFoRCE}~\cite{reforce2025} incorporates self-refinement that effectively retains and reuses execution feedback within a session, but without a structured persistent memory across sessions.
\textsc{Memo-SQL}~\cite{yang2026memosql} maintains a static experience bank of error-repair quintuples constructed offline; while it provides memory, it cannot adapt to novel error patterns at deployment. 
\textsc{Lpe-SQL}~\cite{chu2024lpesql} uses dual notebooks with a Rethink-and-Update mechanism but writes every query outcome immediately, risking noise accumulation. 
\textsc{Mistake Notebook Learning}~\cite{su2025mnl} clusters failures in batches and demonstrates that batch accumulation (batch-16) outperforms per-instance writing on a database QA task.
Our \textsc{SQLSkill} advances this line by (1) employing \emph{deferred batch distillation} with an evidence-gated write (L0--L3 grades and $N_{\text{distill}}$ threshold), ensuring updates rest on statistically sufficient repeated observations; (2) separating episodic memory (Case Library) from procedural memory (Skill Bank) with distinct update frequencies; and (3) tracking errors along four independent semantic dimensions (Table, Field, Relation, Filtering) rather than overlapping syntactic categories. Table~\ref{tab:comparison} summarizes the differences.

\subsection{Error Pattern Distillation}
\label{subsec:distillation}

Extracting reusable error knowledge from failures is critical for agent improvement. \textsc{Reflexion}~\cite{shinn2023reflexion} generates natural-language self-reflections from environment feedback, but these reflections are ephemeral (lost across tasks) and unstructured. \textsc{ExpeL}~\cite{zhao2023expel} persists experiences as free-text rules, but lacks SQL-specific structure: trigger conditions, execution steps, schema applicability, and polarity. \textsc{ReFoRCE}~\cite{reforce2025} generates per-query self-correction guidelines but does not aggregate patterns across queries. \textsc{SQLSkill} contributes \emph{Anti-Pattern Skills}---negative constraints that encode common failure patterns---alongside positive Skills in a unified Skill Bank with conflict resolution routing. This dual-track design (exploratory and governance) enables the agent to learn both what to do and what to avoid. Inspired by \textsc{TextGrad}~\cite{yuksekgonul2024textgrad}'s insight that semantic feedback can be treated as optimizable variables, we replace immediate writes with deferred batch distillation, motivated by the discrete and noisy nature of SQL execution feedback.

\subsection{Agent Memory and Skill Evolution}
\label{subsec:memory}

Long-term memory for LLM agents has been explored in general domains. \textsc{MemGPT}~\cite{packer2023memgpt} manages context windows with tiered storage, while \textsc{Voyager}~\cite{wang2023voyager} stores executable code skills in a skill library. For Text-to-SQL, \textsc{Memo-SQL}~\cite{yang2026memosql} maintains a static experience bank.
\textsc{SkillClaw}~\cite{ma2026skillclaw} proposes a multi-agent skill evolution framework for embodied tasks, with mechanisms for skill merge, split, and deprecation. Our \textsc{SQLSkill} formalizes ``when to write'' via L0--L3 evidence grades and introduces a \emph{Skill Evolver} that performs lifecycle management (merge, split, deprecate) based on counterfactual replay credit assignment.
Unlike prior work, we use interventional evidence (counterfactual gain) rather than observational heuristics to guide skill updates. 
The deferred batch mechanism prevents premature generalization from noisy single-instance feedback. Table~\ref{tab:comparison} provides a structured comparison across the three research gaps identified in Section~\ref{sec:introduction}.

\begin{table*}[t]
  \centering
  \small
  \caption{Comparison of \textsc{SQLSkill} with representative methods across three research gaps. G1: cross-query failure accumulation; G2: SQL-aware structured representation; G3a: quality-gated writing (non-immediate); G3b: positive/negative dual-track Skills.}
  \label{tab:comparison}
  \begin{tabular}{lcccc}
    \toprule
    \textbf{Method} & \textbf{G1: Cross-Query} & \textbf{G2: SQL-Structured} & \textbf{G3a: Write Gate} & \textbf{G3b: $\pm$ Skills} \\
    \midrule
    \textsc{Din-SQL}~\cite{pourreza2023din} & \ding{55} & \ding{55} & \ding{55} & \ding{55} \\
    \textsc{Mac-SQL}~\cite{wang2025macsql} & \ding{55} & \ding{55} & \ding{55} & \ding{55} \\
    \textsc{Chase-SQL}~\cite{chasesql2024} & \ding{55} & \ding{55} & \ding{55} & \ding{55} \\
    \textsc{Reflexion}~\cite{shinn2023reflexion} & \ding{55} & \ding{55} & \ding{55} & \ding{55} \\
    \textsc{ExpeL}~\cite{zhao2023expel} & \ding{51} & \ding{55} & \ding{55} & \ding{55} \\
    \textsc{Memo-SQL}~\cite{yang2026memosql} & \ding{51} (static) & \ding{55} (tuple-based) & \ding{55} (immediate) & \ding{55} \\
    \textsc{ReFoRCE}~\cite{reforce2025} & \ding{51} (self-refine) & \ding{51} (structured agent) & \ding{55} (immediate) & \ding{55} \\
    \midrule
    \textbf{\textsc{SQLSkill} (ours)} & \ding{51} (L0--L3 grades) & \ding{51} (5-tuple + 4 dim.) & \ding{51} ($N_{\text{distill}}$) & \ding{51} (dual-track) \\
    \bottomrule
  \end{tabular}
\end{table*}
